% ==================== Abstract ====================
\begin{abstract}
Real-pair supervised single-image reflection removal (SIRR) trains on images whose
reflection-free target is captured under a protocol, and every protocol silently assumes
something about the reflection layer. We examine those assumptions for \emph{automotive tinted
glass}, where a camera-side occluder (the photographer) cuts part of the reflected solid angle
\emph{and} part of the illumination. Measuring 82 human-confirmed pairs, the two effects are \emph{not} the same size and do not even
point the same way for the two layers: the reflection-free target is only $3.7\%$ darker inside
$\Omega$, while the reflected content is $16.8\%$ weaker, so the input is $13\%$ darker than its
own target there ($I/T|_{\mathrm{in}}\!\approx\!0.87$). A correction applied at one global
strength therefore fails inside the occluded region; the resulting over-correction ratio is
measured at $2.7\times$. Because the occlusion region
cannot be localised from a single frame (IoU ceiling $0.494$ in-domain, $0.14$ across
domains), we do not score methods inside it on the test panel at all

We contribute (i) a tinted-glass benchmark of 631 pairs with four controlled attenuation
levels, per-frame occlusion masks and a signed reflection layer, plus zero-shot baselines for
seven public methods; (ii) a formalization of the occlusion region and a proof that the global
photometric scale is not identifiable from a single frame; (iii) a diagnostic of where
supervision should be placed: with three seeds, an area-matched \emph{random} region
($19.90$~dB) matches the true occlusion mask ($19.41$~dB, difference $-0.49$~dB, not
significant), and a whole-image relative error already reaches $20.05$~dB, so the active
ingredient is the reflection-normalized error rather than the region identity, while a
closed-form support-area bound ($0.036$~dB) shows why boundary-scale terms cannot be validated
at all; and (iv) a transfer study: fine-tuning on our data gains $\mathbf{+2.22}$~dB over the
best public method on \textbf{15/15} scenes and still $+1.69$~dB after a per-image affine, and
the gain is largest for the pool that contains \emph{the fewest} occlusion regions --- what
transfers is the domain, not the occlusion. An interventional protocol shows that
state-of-the-art photometric conditioning fails its own requirement ($0.02$ vs.\ $0.50$).
\end{abstract}

\begin{IEEEkeywords}
single-image reflection removal, tinted glass, occlusion region, identifiability,
supervision placement, benchmark
\end{IEEEkeywords}
